\chapter{GIỚI THIỆU ĐƠN VỊ THỰC TẬP}

\section{Thông tin chung}

\CompanyName\ (tên quốc tế: \textit{TECHPAL Joint Stock Company}, viết tắt
TECHPAL., JSC) có mã số thuế 0104469357, thành lập từ ngày 08/02/2010
\cite{techpal-intro}. Địa chỉ
văn phòng giao dịch đặt tại \CompanyAddress. Số điện thoại liên hệ là 024 3793
8272, email \href{mailto:info@techpal.vn}{info@techpal.vn}, website
\url{https://techpal.vn}.

\begin{table}[H]
\centering
\caption{Thông tin khái quát về đơn vị thực tập}
\label{tab:company-info}
\begin{tabularx}{\textwidth}{>{\bfseries}p{4.3cm}X}
\toprule
Tên doanh nghiệp & Công ty Cổ phần TECHPAL \\
Tên quốc tế & TECHPAL Joint Stock Company \\
Tên viết tắt & TECHPAL., JSC \\
Mã số thuế & 0104469357 \\
Ngày thành lập & 08/02/2010 \\
Trụ sở chính & \CompanyAddress \\
Website & \url{https://techpal.vn} \\
Người hướng dẫn & \CompanySupervisor \\
Đại diện xác nhận & \CompanyRepresentative, Phó Tổng giám đốc \\
Thời gian thực tập & \InternshipPeriod \\
\bottomrule
\end{tabularx}
\end{table}

\section{Quá trình hình thành và định hướng phát triển}

Trên website chính thức, TECHPAL giới thiệu bốn mảng hoạt động chính gồm: kinh
doanh linh kiện điện tử và thiết bị tự động hóa; nghiên cứu, chuyển giao công
nghệ; năng lượng xanh (chủ yếu là điện mặt trời); và đầu tư nông nghiệp công
nghệ cao \cite{techpal-intro}. Nhìn chung, công ty xuất phát từ lĩnh vực điện --
điện tử rồi mở rộng thêm sang năng lượng và nông nghiệp.

Về định hướng, TECHPAL tập trung vào đổi mới công nghệ và phát triển bền vững.
Công ty đã triển khai một số dự án điện mặt trời áp mái, song song với hoạt động
cung ứng thiết bị tự động hóa và các dự án nông nghiệp ứng dụng công nghệ
\cite{techpal-intro}.

\section{Lĩnh vực hoạt động chuyên môn}

\subsection{Linh kiện điện tử và thiết bị tự động hóa}

Trên trang bán hàng của TECHPAL, sản phẩm được chia thành bốn nhóm: linh kiện
điện tử, thiết bị tự động hóa công nghiệp, vật tư phụ trợ và thiết bị phòng sạch.
Em thấy trên danh mục có khá nhiều mã linh kiện bán dẫn, đầu nối, rơ-le, nguồn
cấp, cảm biến, bộ điều khiển PLC, mô-đun vào/ra và bộ ghi dữ liệu
\cite{techpal-components}.

Việc cung ứng linh kiện kỹ thuật không đơn giản chỉ là bán hàng. Người làm kỹ
thuật phải đọc đúng mã linh kiện, so sánh thông số điện và cơ khí, kiểm tra xem
có thay thế được không, xem sản phẩm còn được sản xuất hay đã ngừng, rồi cả
điều kiện lưu kho, vận chuyển và chống tĩnh điện nữa. Bỏ sót bước nào cũng có
thể gây ra lỗi cho hệ thống về sau.

\subsection{Nghiên cứu và chuyển giao công nghệ}

Mảng R\&D của TECHPAL hoạt động như cầu nối giữa yêu cầu thực tế và giải pháp
kỹ thuật. Quy trình thường bắt đầu từ khảo sát nhu cầu, sau đó phân rã chức
năng, chọn phần cứng, thiết kế kết nối, thử nghiệm rồi bàn giao. Với sinh viên
ngành Điện tử Viễn thông, đây là nơi có thể áp dụng các kiến thức về mạch điện,
vi điều khiển, cảm biến hay truyền thông công nghiệp.

\subsection{Năng lượng xanh và điện mặt trời}

Một hệ thống điện mặt trời thông thường gồm tấm pin, thiết bị bảo vệ DC, bộ
biến đổi (inverter), hệ thống giám sát và phần kết nối lưới. Khi chọn thiết bị
cho hệ thống này cần xét đến công suất, dải điện áp, hiệu suất, điều kiện thời
tiết và yêu cầu an toàn. TECHPAL có tham gia lĩnh vực này thông qua các dự án
điện mặt trời áp mái.

\subsection{Nông nghiệp công nghệ cao}

Mảng nông nghiệp công nghệ cao tạo ra nhu cầu về cảm biến môi trường, bộ điều
khiển, thiết bị truyền thông và lưu trữ dữ liệu. Tuy mục tiêu cuối cùng là phục
vụ sản xuất nông nghiệp, nhưng nhiều thiết bị cốt lõi vẫn thuộc lĩnh vực điện
tử và tự động hóa -- ví dụ như đo nhiệt độ, độ ẩm, điều khiển quạt, bơm hay thu
thập dữ liệu từ xa.

\section{Cơ cấu tổ chức}

Theo thông tin trên website, TECHPAL có cơ cấu gồm Đại hội đồng cổ đông, Hội
đồng quản trị, Ban Kiểm soát, Ban Tổng giám đốc và các phòng chuyên môn
\cite{techpal-intro}. Cụ thể,
các phòng ban được liệt kê gồm: Kinh doanh linh kiện điện tử, Tài chính -- Kế
toán, Hành chính -- Tổng hợp, Phát triển và vận hành dự án, Pháp chế doanh
nghiệp. Ngoài ra còn có đơn vị thành viên, văn phòng đại diện và các ban quản lý
dự án. Bài đăng về Đại hội đồng cổ đông thường niên năm
2026 cũng xác nhận ông \CompanyRepresentative\ giữ chức Thành viên Hội đồng quản
trị kiêm Phó Tổng giám đốc \cite{techpal-agm2026}.

\section{Các vị trí công việc và yêu cầu năng lực}

Từ các lĩnh vực hoạt động nêu trên, em tạm chia ra một số nhóm vị trí liên quan
đến chuyên ngành Điện tử Viễn thông như trong Bảng~\ref{tab:positions}. Cần lưu ý
đây là cách em tự phân nhóm dựa trên tính chất công việc, không phải danh sách
tuyển dụng chính thức của công ty.

\begin{table}[htbp]
\centering
\caption{Một số nhóm vị trí và năng lực cần thiết}
\label{tab:positions}
\small
\renewcommand{\arraystretch}{1.08}
\begin{tabularx}{\textwidth}{>{\raggedright\arraybackslash}p{3.1cm}
                  >{\raggedright\arraybackslash}X
                  >{\raggedright\arraybackslash}X}
\toprule
\textbf{Nhóm vị trí} & \textbf{Nhiệm vụ điển hình} & \textbf{Kiến thức, kỹ năng cần có} \\
\midrule
Kỹ sư ứng dụng linh kiện & Đọc yêu cầu, chọn linh kiện, kiểm tra khả năng thay thế và hỗ trợ kỹ thuật & Mạch điện tử, đọc datasheet, tra cứu mã hàng, giao tiếp kỹ thuật \\
\addlinespace
Kỹ sư tự động hóa & Lựa chọn PLC, cảm biến, cơ cấu chấp hành; cấu hình và kiểm thử hệ thống & Điều khiển, điện công nghiệp, PLC/HMI, truyền thông công nghiệp \\
\addlinespace
Kỹ sư dự án & Khảo sát, lập phương án, bóc tách thiết bị, theo dõi triển khai & Tư duy hệ thống, đọc bản vẽ, lập kế hoạch, quản lý rủi ro \\
\addlinespace
Kỹ sư R\&D & Phân tích yêu cầu, thiết kế nguyên mẫu, đo kiểm và lập tài liệu & Điện tử tương tự -- số, vi điều khiển, đo lường, thiết kế và kiểm thử \\
\addlinespace
Kiểm tra kỹ thuật & Kiểm tra mã hàng, ngoại quan, thông số và hồ sơ chất lượng & Dụng cụ đo, tiêu chuẩn chất lượng, ESD, kỹ năng ghi nhận dữ liệu \\
\bottomrule
\end{tabularx}
\end{table}
